%==============================================================================
% jarvis-listings.tex
% Code-listing setup for the Jarvis-HEP manual.
%
% Provides:
%   - a Jarvis color palette
%   - a custom YAML language definition (listings has none built in)
%   - styles:  yaml   (task-card YAML)
%              bash   (shell / CLI sessions)
%              plain  (generic / SLHA / output dumps)
%   - convenience inline/float macros
%
% Long real-world cards should be \lstinputlisting'd from ../listings/ so they
% stay byte-true with jarvishep/project_template.  Inside Tufte's narrow text
% column, wrap wide listings in a {fullwidth} environment.
%==============================================================================

\usepackage{listings}

% --- Palette ------------------------------------------------------------------
\definecolor{jBg}{HTML}{F7F7F4}      % listing background
\definecolor{jRule}{HTML}{D8D8D0}    % frame rule
\definecolor{jKey}{HTML}{1F4E79}     % YAML keys / keywords (deep blue)
\definecolor{jStr}{HTML}{8A5A00}     % strings (amber)
\definecolor{jComment}{HTML}{6E7B6E} % comments (muted green-grey)
\definecolor{jNum}{HTML}{7A1F8B}     % numbers (purple)
\definecolor{jMarker}{HTML}{B3261E}  % &J/ @tokens / ${...} (Jarvis red)
\definecolor{jLineNo}{HTML}{A0A0A0}
\definecolor{jOut}{HTML}{4A5568}      % what a program printed back: muted slate,
                                      % one step lighter than jShell (334155) so the
                                      % output half of a `terminal' box reads as a
                                      % reply rather than as something you type
\definecolor{jLineNum}{HTML}{C7CCD2}  % very light grey for YAML line numbers
\definecolor{jGuide}{HTML}{E3E6EA}    % very light grey for YAML indentation guides

% --- Semantic colours (used inline and in listings) ---------------------------
\definecolor{jExpr}{HTML}{C2185B}     % expressions  -> pink
\definecolor{jPath}{HTML}{1B5E20}     % file/dir paths -> dark green
\definecolor{jCmd}{HTML}{E8730C}      % Jarvis CLI commands -> Jarvis orange
\definecolor{jVar}{HTML}{C2185B}      % variables / arguments -> pink (same hue as jExpr)
\definecolor{jFunc}{HTML}{1B5E20}     % function / operator names -> dark green (same hue as jPath)
\definecolor{jYKey}{HTML}{1971C2}     % YAML keys, in listings and in prose (bright blue)
\definecolor{jYValue}{HTML}{2F9E44}   % YAML (string) values, in listings and in prose (bright green)
\definecolor{jYBool}{HTML}{A6462E}    % YAML true/false/null, in listings
\definecolor{jSampler}{HTML}{7A1F8B}  % sampler method names (same purple family as jNum)

% --- Inline semantic markup ---------------------------------------------------
\newcommand{\expr}[1]{\textcolor{jExpr}{\jtt{#1}}}   % an expression, e.g. \expr{x * Pi}
\newcommand{\cmd}[1]{\textcolor{jCmd}{\jtt{#1}}}     % a Jarvis CLI command, e.g. \cmd{Jarvis}
% \variable{} / \mvar{}: a variable or argument name -> pink.
% \variable{} is for monospace/code contexts (e.g. Sampling.Variables[].name);
% \mvar{} is for math mode, e.g. $\mvar{x}$, $f(\mvar{x},\mvar{y})$.
\newcommand{\variable}[1]{\textcolor{jVar}{\jtt{#1}}}
\newcommand{\mvar}[1]{\textcolor{jVar}{#1}}
% \func{} / \mfunc{}: a function or operator name -> dark green.
% \func{} is for monospace/code contexts, e.g. \func{sin}; \mfunc{} is for math
% mode operator names, e.g. $\mfunc{\sin}(\mvar{x})$.
\newcommand{\func}[1]{\textcolor{jFunc}{\jtt{#1}}}
\newcommand{\mfunc}[1]{\textcolor{jFunc}{#1}}
% \ykey{} / \yvalue{}: a YAML key / value mentioned in prose, coloured to match
% their appearance inside a yamlcode listing (jYKey blue, jYValue green) --
% \yamlkey{} (jarvis-preamble.tex) is the same key colour, kept as the
% established name; \ykey{} is a shorter alias for new prose.
\newcommand{\ykey}[1]{\textcolor{jYKey}{\jtt{#1}}}
\newcommand{\yvalue}[1]{\textcolor{jYValue}{\jtt{#1}}}

% --- A minimal YAML dialect ---------------------------------------------------
% Good enough to colour keys, strings, comments, booleans, and Jarvis markers
% without a full parser.
\lstdefinelanguage{jyaml}{
  morekeywords={true,false,null,True,False,None,yes,no},
  sensitive=true,
  morecomment=[l]{\#},
  morestring=[b]",
  morestring=[b]',
  alsoletter={-,_,.},
}

% --- Shared base style --------------------------------------------------------
% `listings' processes verbatim material at the catcode level, before UTF-8
% decoding -- a literal non-ASCII byte inside any lstlisting-based environment
% (yamlcode/bashcode/plaincode) is a hard "Invalid UTF-8 byte sequence" error,
% not a warning. `literate' maps the specific bytes real console/log output
% actually contains (the Jarvis-logger bullet, plus-minus in evidence lines)
% to their LaTeX-safe renderings, so real excerpts can be quoted verbatim.
% Add more pairs here if a future quote needs another character; do NOT type
% the raw byte directly into a listing without a matching entry.
\lstdefinestyle{jbase}{
  backgroundcolor=\color{jBg},
  basicstyle=\ttfamily\footnotesize,
  breaklines=true,
  breakatwhitespace=false,
  postbreak=\mbox{\textcolor{jLineNo}{$\hookrightarrow$}\space},
  columns=fullflexible,
  keepspaces=true,
  showstringspaces=false,
  tabsize=2,
  frame=leftline,
  framerule=1.2pt,
  rulecolor=\color{jRule},
  xleftmargin=0.8em,
  framexleftmargin=0.8em,
  aboveskip=0.9\baselineskip,
  belowskip=0.6\baselineskip,
  literate=
    {±}{{$\pm$}}1
    {·}{{\textperiodcentered}}1
    {•}{{\textbullet}}1,
  % NOTE on box-drawing bytes (╭ │ ─ ━ ╰ …): the CLI renders its help and runtime
  % panels with Rich, and quoting those frames verbatim was tried and REVERTED.
  % pmboxdraw's glyphs do not advance exactly one monospace cell, so `literate'
  % width hints cannot hold the columns and every panel line wrapped. Reproduce
  % such output with the frame stripped instead -- the `terminal' box already
  % supplies a border, and Rich's own frame width is a property of the reader's
  % terminal, not of the program.
}

% --- YAML style ---------------------------------------------------------------
% Keys (and other bare/unquoted tokens) take the default text colour, jYKey;
% quoted string values take jYValue; true/false/null take jYBool -- the same
% three colours \ykey{}/\yvalue{} use inline, so a key or value reads the same
% whether it is quoted in prose or sitting inside the listing box.
\lstdefinestyle{yaml}{
  style=jbase,
  language=jyaml,
  basicstyle=\ttfamily\footnotesize\color{jYKey},
  keywordstyle=\color{jYBool}\bfseries,  % booleans/null
  stringstyle=\color{jYValue},
  commentstyle=\color{jComment}\itshape,
  % Highlight Jarvis markers and tokens.
  moredelim=[is][\color{jMarker}\bfseries]{@@}{@@},
  % Very light line numbers in the left margin (outside the frame rule).
  numbers=left,
  numberstyle=\tiny\color{jLineNum},
  numbersep=9pt,
  xleftmargin=2.6em,
  framexleftmargin=2.6em,
}

% --- Bash / CLI style ---------------------------------------------------------
\lstdefinestyle{bash}{
  style=jbase,
  language=bash,
  keywordstyle=\color{jKey}\bfseries,
  commentstyle=\color{jComment}\itshape,
  stringstyle=\color{jStr},
  morekeywords={Jarvis,pip,python3,cd,latexmk},
}

% --- Plain dump style (SLHA blocks, output text) ------------------------------
\lstdefinestyle{plain}{
  style=jbase,
  language={},
  keywordstyle=\color{jKey},
  commentstyle=\color{jComment}\itshape,
}

% --- Inner style for the boxed shell environment (frame/bg come from the box) --
\lstdefinestyle{shellinner}{
  language=bash,
  % Self-contained on purpose. The document-level \lstset{style=jbase} is ambient,
  % so any decoration jbase sets and this style does not override leaks in -- which
  % is how the `terminal' box first grew a stray grey left tick and a jBg fill that
  % painted over the box's own (differing) upper/lower background tints.
  frame=none, backgroundcolor={}, numbers=none,
  xleftmargin=0pt, framexleftmargin=0pt,
  basicstyle=\ttfamily\footnotesize,
  keywordstyle=\color{jKey}\bfseries,
  commentstyle=\color{jComment}\itshape,
  stringstyle=\color{jStr},
  morekeywords={pip,python3,cd,source,pyenv,latexmk,make},
  keywords=[2]{Jarvis,jplot,jopera},        % Jarvis-family executables
  keywordstyle=[2]\color{jCmd}\bfseries,     % ... in Jarvis orange
  breaklines=true,
  postbreak=\mbox{\textcolor{jLineNo}{$\hookrightarrow$}\space},
  columns=fullflexible, keepspaces=true, showstringspaces=false, tabsize=2,
  aboveskip=0pt, belowskip=0pt,   % no extra listing space inside the box
  literate=
    {±}{{$\pm$}}1
    {·}{{\textperiodcentered}}1
    {•}{{\textbullet}}1,
  % NOTE on box-drawing bytes (╭ │ ─ ━ ╰ …): the CLI renders its help and runtime
  % panels with Rich, and quoting those frames verbatim was tried and REVERTED.
  % pmboxdraw's glyphs do not advance exactly one monospace cell, so `literate'
  % width hints cannot hold the columns and every panel line wrapped. Reproduce
  % such output with the frame stripped instead -- the `terminal' box already
  % supplies a border, and Rich's own frame width is a property of the reader's
  % terminal, not of the program.
}

% --- Inner style for the OUTPUT half of the boxed terminal environment ---------
% Deliberately quieter than `shellinner': what a program printed back is not
% something you type, so nothing here is highlighted as a command. Only the
% Jarvis-family names keep their colour (they identify the component that spoke),
% and log levels stay legible. Same metrics as `shellinner' so the two halves of
% one `terminal' box line up character for character across the divider.
\lstdefinestyle{termout}{
  language={},
  frame=none, backgroundcolor={}, numbers=none,   % see the note in `shellinner'
  xleftmargin=0pt, framexleftmargin=0pt,
  basicstyle=\ttfamily\footnotesize\color{jOut},
  keywords=[2]{Jarvis,Jarvis2,jplot,jopera},
  keywordstyle=[2]\color{jCmd},
  breaklines=true,
  postbreak=\mbox{\textcolor{jLineNo}{$\hookrightarrow$}\space},
  columns=fullflexible, keepspaces=true, showstringspaces=false, tabsize=2,
  aboveskip=0pt, belowskip=0pt,
  literate=
    {±}{{$\pm$}}1
    {·}{{\textperiodcentered}}1
    {•}{{\textbullet}}1,
  % NOTE on box-drawing bytes (╭ │ ─ ━ ╰ …): the CLI renders its help and runtime
  % panels with Rich, and quoting those frames verbatim was tried and REVERTED.
  % pmboxdraw's glyphs do not advance exactly one monospace cell, so `literate'
  % width hints cannot hold the columns and every panel line wrapped. Reproduce
  % such output with the frame stripped instead -- the `terminal' box already
  % supplies a border, and Rich's own frame width is a property of the reader's
  % terminal, not of the program.
}

% --- YAML style for the "window" box (frame/background come from the box) ------
% Same key/value/boolean colours as the plain `yaml' style (jYKey/jYValue/jYBool);
% additionally, expression values (after `expression: "') go pink and &J paths
% go dark green -- the window box is the one place both distinctions show at once.
\lstdefinestyle{yamlwindow}{
  language=jyaml,
  % No colour here: the enclosing `yamlwin' box supplies it via `colupper' so
  % it survives a page-break split (see jarvis-boxes.tex for why).
  basicstyle=\ttfamily\footnotesize,
  keywordstyle=\color{jYBool}\bfseries,
  commentstyle=\color{jComment}\itshape,
  morestring=[b]", stringstyle=\color{jYValue},
  moredelim=[s][\color{jExpr}]{expression: "}{"},
  moredelim=[l][\color{jPath}]{"\&J},
  numbers=left, numberstyle=\tiny\color{jLineNum}, numbersep=12pt,
  xleftmargin=2.4em,
  breaklines=true, columns=fullflexible, keepspaces=true, showstringspaces=false, tabsize=2,
  aboveskip=0pt, belowskip=0pt,
  literate=
    {±}{{$\pm$}}1
    {·}{{\textperiodcentered}}1
    {•}{{\textbullet}}1,
  % NOTE on box-drawing bytes (╭ │ ─ ━ ╰ …): the CLI renders its help and runtime
  % panels with Rich, and quoting those frames verbatim was tried and REVERTED.
  % pmboxdraw's glyphs do not advance exactly one monospace cell, so `literate'
  % width hints cannot hold the columns and every panel line wrapped. Reproduce
  % such output with the frame stripped instead -- the `terminal' box already
  % supplies a border, and Rich's own frame width is a property of the reader's
  % terminal, not of the program.
}

\lstset{style=jbase}

% --- Convenience environments / macros ----------------------------------------
% \begin{yamlcode} ... \end{yamlcode}
\lstnewenvironment{yamlcode}{\lstset{style=yaml}}{}
% \begin{bashcode} ... \end{bashcode}
\lstnewenvironment{bashcode}{\lstset{style=bash}}{}
% \begin{plaincode} ... \end{plaincode}
\lstnewenvironment{plaincode}{\lstset{style=plain}}{}

% File include shortcuts (keep wide ones inside {fullwidth}).
\newcommand{\yamlfile}[1]{\lstinputlisting[style=yaml]{#1}}
\newcommand{\bashfile}[1]{\lstinputlisting[style=bash]{#1}}

% Inline monospace shortcut already provided as \code in the preamble.
